\documentclass[11pt]{article}
\usepackage[margin=1in]{geometry}
\usepackage{amsmath,amssymb,amsthm,mathtools,microtype,hyperref}
\hypersetup{hidelinks}
\newtheorem{theorem}{Theorem}
\newtheorem{lemma}[theorem]{Lemma}
\newcommand{\ddc}{dd^c}
\newcommand{\dc}{d^c}
\newcommand{\dA}{\mathrm dA}
\newcommand{\CP}{\mathbb P}
\newcommand{\C}{\mathbb C}
\title{A nef and big divisorial equation with an incomplete null end}
\author{}
\date{4 October 2026}
\begin{document}
\maketitle
\begin{abstract}
We construct a compact K\"ahler surface, a nef and big background and a
divisorial exponential Monge--Amp\`ere equation with a globally smooth
strictly positive pole coefficient. Its solution has full mass, finite
energy and zero Lelong numbers, but its positive metric is incomplete on
every dense Zariski open regular locus. The failure occurs at a null
exceptional curve disjoint from the pole divisor. Two exact radial ends
and a positive-derivative interpolation give all the data. Exact Jacobian
cancellation proves smoothness of the volume, and radial shell masses
prove finite energy by bounded truncations. Mere nefness and bigness
therefore do not suffice for completeness; additional geometric-smoothing
compatibility is a separate question.
\end{abstract}

\section{The completeness question}
Use $\ddc=(i/(2\pi))\partial\bar\partial$ and Euclidean area $\dA$.
For a compact K\"ahler manifold, a big class represented by $\theta$,
a simple-normal-crossing divisor $D$ with smooth squared section norm
$|\sigma|^2$, and a smooth strictly positive volume $dV_Y$, consider
\begin{equation}\label{eq:equation}
 \langle(\theta+\ddc\rho)^n\rangle=e^\rho|\sigma|^{-2}dV_Y.
\end{equation}
The product is non-pluripolar and the density is extended by zero on $D$.
The equation's full-mass finite-energy solution is the one in the
completeness question of \cite[Section 3.1]{DGL}. We ask whether nefness
and bigness alone force a complete positive metric on some dense Zariski
open set where that current is smooth and strictly positive.

\begin{theorem}\label{thm:main}
There are such data on a compact connected K\"ahler surface $Y$, with
$\theta$ smooth closed semipositive and $\{\theta\}$ nef and big, for
which a full-mass solution $\rho\in\mathcal E^1(Y,\theta)$ has zero
Lelong numbers and is smooth on $Y\setminus(E\cup D)$. Its current
$\Omega=\theta+\ddc\rho$ is strictly positive there. The divisors
$E$ and $D$ are disjoint, and $E$ is a null curve. Every dense Zariski
open positive regular locus is incomplete, although every approach to
$D$ within $Y\setminus(E\cup D)$ has infinite length.
\end{theorem}

The prescribed measure has the exact nonintegrable pole
$|\sigma|^{-2}dV_Y$ with a smooth strictly positive coefficient across
$D$. The conclusion concerns mere nefness and bigness. The class,
discrepancy, residual and volume identities in the flow setup of
\cite{DGL} will not be asserted for this example.

\section{An exact radial profile}
Primes in this section denote differentiation in $t\in\mathbb R$.
\begin{lemma}\label{lem:profile}
There is $F\in C^\infty(\mathbb R)$ with $0<p:=F'<1$, $p'>0$,
and constants $c=1/100$, $B\in\mathbb R$, $K>0$, such that
\begin{align}
 F(t)&=B-3\log(1-ce^{t/2}),&
 pp'&=Ke^{F+t},&& t\leq-2,\label{eq:left}\\
 F(t)&=t+\log H(p(t)),&
 pp'&=e^{F-t},&& t\geq2,\label{eq:right}\\
 H(p)&=\frac{(1-p)^2(1+2p)}6,&&&0<p<1.
\end{align}
Moreover,
\begin{gather}
 F(t)=B+3ce^{t/2}+O(e^t),\quad
 p'(t)\sim(3c/4)e^{t/2}\quad(t\to-\infty),\label{eq:left-asymptotic}\\
 1-p(t)\sim2/t,\qquad p'(t)\sim2/t^2\quad(t\to\infty),\notag\\
 F(t)-t=-2\log t+\log2+o(1)\quad(t\to\infty).\label{eq:right-asymptotic}
\end{gather}
\end{lemma}

\begin{proof}
Put $L(t)=-3\log(1-ce^{t/2})$ on a neighborhood of $(-\infty,-2]$.
With $q=ce^{t/2}$, differentiation gives
\[
 p_L=L'=\frac{3q}{2(1-q)},\quad p_L'=\frac{3q}{4(1-q)^2},\quad
 p_Lp_L'=\frac{9c^2}8e^t(1-ce^{t/2})^{-3}.
\]
Thus \eqref{eq:left} holds for any additive $B$, with
$K=(9c^2/8)e^{-B}$; $p_L(-2)<1/50$.

For the other end $H'(p)=(p-1)p$. Define
\[
 \mathcal T(p)=2+\int_{3/4}^p\frac{q}{H(q)}\,dq.
\]
It is strictly increasing, and tends to infinity as $p\uparrow1$,
since its integrand is asymptotic to $2/(1-p)^2$. Its inverse $p_R(t)$
exists for $t\geq2$, extends across $2$, and satisfies
$p_R(2)=3/4$, $p_R'=H(p_R)/p_R>0$. Set
$F_R=t+\log H(p_R)$. The derivative identity for $H$ gives
$F_R'=p_R$, proving \eqref{eq:right}.
Writing $q=1-p_R$ gives
\[
 (1/q)'=\frac{3-2q}{6(1-q)}\longrightarrow\frac12.
\]
Hence $q\sim2/t$, $H(p_R)\sim2/t^2$ and
\eqref{eq:right-asymptotic} follows. This proves existence and all
needed asymptotics without a shooting parameter.

To patch derivatives on $[-2,2]$, let $d=p_R(2)-p_L(-2)>0$.
Choose disjoint short endpoint collars and cutoffs $\alpha_L,\alpha_R$
equal to one on smaller collars and supported in the larger ones. For
a small $\varepsilon>0$ the function
\[
 Q_0=\alpha_Lp_L'+\alpha_Rp_R'
              +\varepsilon(1-\alpha_L-\alpha_R)
\]
is strictly positive and matches both endpoint derivative germs.
Shortening the collars and decreasing $\varepsilon$ makes its integral
less than $d$. Add a nonnegative smooth interior bump with integral
$d-\int Q_0$ to obtain $Q>0$ with integral $d$.
Then $p=p_L(-2)+\int_{-2}^tQ$ agrees with both endpoint functions
on their smaller collars, so it patches smoothly to $p_L,p_R$.
It is increasing and stays in $(0,1)$. Choose $F(2)=F_R(2)$ and
integrate $F'=p$. The left primitive differs from $L$ by a constant $B$;
the right one equals $F_R$. The left expansion is immediate.
\end{proof}

\section{The compact data}
Let $\pi:Y\to\CP^2$ be the blowup at $[1:0:0]$, represented by
\[
 Y=\{([Z_0:Z_1:Z_2],[\ell_1:\ell_2])\in\CP^2\times\CP^1:
                                      Z_1\ell_2=Z_2\ell_1\}.
\]
This smooth incidence surface is compact, connected and K\"ahler.
Write $f$ for its second projection, $E$ for the exceptional curve,
and $D=\pi^{-1}(Z_0=0)$. The latter is a smooth projective line
disjoint from $E$. Let $H=\pi^*c_1(\mathcal O_{\CP^2}(1))$.
It has a smooth semipositive representative, so it is nef, and $H^2=1$.
For bigness, $f^*c_1(\mathcal O_{\CP^1}(1))=H-[E]$, as seen from
the total transform of a line through the blown-up point. If $\beta,\eta$
are the two pulled-back Fubini--Study forms, then $\beta+\varepsilon\eta$
is K\"ahler, by the incidence embedding, and
\[
 \frac{\beta+\varepsilon\eta}{1+\varepsilon}
                         +\frac\varepsilon{1+\varepsilon}[E]
\]
is a K\"ahler current in $H$. This proves bigness directly.

Choose an even nonnegative smooth function $\zeta$, positive on $(-1,1)$,
supported in $[-1,1]$, of integral one. Define
\[
 U(t)=\int_{\mathbb R}\max(0,t-x)\zeta(x)\,dx.
\]
Then $U$ is smooth convex, $U=0$ for $t\leq-1$, $U=t$ for
$t\geq1$, and $U',U''\geq0$. On $Z_0\ne0$, set
$z=(Z_1/Z_0,Z_2/Z_0)$, $s=|z|^2$, $t=\log s$.
Give the frame $Z_0$ of $\mathcal O(1)$ squared norm $e^{-U(t)}$.
Near the origin this is constant. At infinity, in coordinates
\[
 w=Z_0/Z_1,\qquad v=Z_2/Z_1,\qquad h=1+|v|^2,
 \qquad s=h/|w|^2,
\]
the frame $Z_1$ has norm $e^{-U}/|w|^2=1/h$ near $w=0$.
Thus this is a smooth strictly positive bundle metric globally.
Pull it back to $\mathcal O_Y(D)$ and write $\sigma=\pi^*Z_0$.
Its curvature $\theta$ is smooth closed semipositive in $H$, with
\begin{align}
 |\sigma|^2&=e^{-U(t)},&\theta&=\ddc U(t)&&\text{on the affine chart},\label{eq:background}\\
 |\sigma|^2&=|w|^2/h,&\theta&=\ddc\log h&&\text{near infinity}.
\end{align}
Semipositivity follows from the radial Hessian and $U',U''\geq0$.
The form vanishes near $E$, and has rank one near $D$.
The normalized minimal potential is zero; the volume of the class is
$\int_Y\theta^2=H^2=1$, as in the semipositive volume description
of \cite[Sections 1.5 and 4]{BEGZ}.

\section{The solution and exact smooth volume}
On $Y\setminus(E\cup D)\simeq\C^2\setminus\{0\}$ put
\begin{equation}\label{eq:solution}
 \rho=F(t)-U(t),\qquad \Omega=\ddc F(t),\qquad
 dV_Y=e^{-F(t)}\Omega^2.
\end{equation}
The radial Hessian of $F(\log|z|^2)$ has eigenvalues $p/s,p'/s$.
Both are positive, and with our convention
\begin{equation}\label{eq:determinant}
 \Omega^2=\frac2{\pi^2}\frac{pp'}{s^2}\dA_{z_1}\dA_{z_2}.
\end{equation}
Extend $\rho$ to $E$ by its limit $B$ and to $D$ by $-\infty$.
Near $E$, use $z=(u,uv)$, so $s=|u|^2h$. Its local potential is
\[
 F(t)=B-3\log(1-c|u|\sqrt h).
\]
It is the pullback of a locally bounded plurisubharmonic extension at
the affine origin, so it is continuous and plurisubharmonic near $E$.
Near $D$, the local potential is $\rho+\log h$, which is psh off
$w=0$, locally bounded above and tends to $-\infty$.
Its bounded-above psh extension \cite[Proposition 3.3.25]{JP} proves that the global $\rho$ is
$\theta$-psh. Boundedness at $E$ and the log-log growth at $D$ give
zero Lelong numbers everywhere.

We verify smoothness of the volume through exact identities. At $E$
the Jacobian is $|u|^2$, so
\[
 \Omega^2=\frac2{\pi^2}\frac{pp'}{|u|^2h^2}\dA_u\dA_v.
\]
Using \eqref{eq:left}, $e^{-F}pp'=Ks=K|u|^2h$, and therefore
\begin{equation}\label{eq:volume-E}
 dV_Y=\frac{2K}{\pi^2h}\dA_u\dA_v
 \qquad\text{near }E.
\end{equation}
This extends smoothly and strictly positively across $u=0$.
At $D$, $z=(1/w,v/w)$ has real Jacobian $|w|^{-6}$, so
\[
 \Omega^2=\frac2{\pi^2}\frac{pp'}{|w|^2h^2}\dA_w\dA_v.
\]
Now \eqref{eq:right} gives $e^{-F}pp'=e^{-t}=|w|^2/h$ and
\begin{equation}\label{eq:volume-D}
 dV_Y=\frac2{\pi^2h^3}\dA_w\dA_v
 \qquad\text{near }D.
\end{equation}
The other exceptional and infinity charts give the same conclusion.
The extensions agree on overlaps since they agree on the dense punctured
overlaps. Between the ends \eqref{eq:solution} is smooth and positive.
Thus $dV_Y$ is globally smooth and strictly positive. In particular,
\[
 |\sigma|^{-2}dV_Y=\frac2{\pi^2|w|^2h^2}\dA_w\dA_v
 \qquad\text{near }D;
\]
its transverse integral diverges, with a smooth positive coefficient.

On the complement the equation is exact:
\[
 e^\rho|\sigma|^{-2}dV_Y
          =e^{F-U}e^Ue^{-F}\Omega^2=\Omega^2.
\]
The non-pluripolar product has no mass on the two divisors, and the
right side has no mass there, with the stipulated extension at $D$.
Consequently the global measure identity holds once its finite mass
is checked below.

\section{Full mass and finite energy by bounded truncations}
Polar integration in $\C^2$ has radial density $\pi^2s\,ds$.
Equation~\eqref{eq:determinant} therefore yields
\begin{equation}\label{eq:mass}
 \int_Y\langle\Omega^2\rangle
                 =\int_{-\infty}^{\infty}2pp'\,dt=1.
\end{equation}
This is the class volume, so the solution has full mass.
Near $E$, $\rho$ is bounded and $\theta=0$. Away from the two ends
all data are smooth on a compact set. On the cusp end write
\[
 \alpha=\ddc t=\ddc\log h,\qquad
 \beta_t=\frac i{2\pi}\partial t\wedge\bar\partial t.
\]
Both have rank one, and
\begin{gather*}
 \theta=\alpha,\qquad \Omega=p\alpha+p'\beta_t,\\
 \Omega^2=2pp'\alpha\wedge\beta_t,\qquad
 \Omega\wedge\theta=p'\alpha\wedge\beta_t,\qquad\theta^2=0.
\end{gather*}
Integrating $\alpha\wedge\beta_t$ in its angular and projective
variables gives $dt$. Since $|\rho|=O(\log t)$ and $p'=O(t^{-2})$,
both $\int|\rho|\Omega^2$ and $\int|\rho|\Omega\wedge\theta$
are finite at that end.

We spell out the bounded approximation to verify the finite-energy
definition. Choose $M$ with $\phi=\rho-M\leq0$, and let
$\phi_k=\max(\phi,-k)$. For large $k$ truncation changes the potential
only near $D$. There $\phi=\log H(p(t))-M$ decreases strictly to
$-\infty$; denote its level $-k$ by $t=T_k$. Then
\[
 k=2\log T_k+O(1).
\]
On the punctured chart the radial primitive for the truncated current is
\[
 G_k(t)=\max(F(t)-M,t-k).
\]
A genuine local psh potential is $G_k(t)+\log|w|^2$, which equals
$\log h-k$ near $D$. Thus no divisor mass is added by writing the
primitive in the punctured coordinate. Its derivative is $p$ before
$T_k$ and $1$ after $T_k$, with an upward jump.
Convolution in $t$ with an even nonnegative smooth kernel gives smooth
convex primitives decreasing to $G_k$ as its scale decreases. The bounded
local potentials converge locally uniformly, and the Bedford--Taylor
products converge \cite[Theorems 2.1 and 2.10]{BT}. Their radial
top-degree measures are $d((G_k')^2)$ and $dG_k'$ for the square and the
wedge with $\theta$. Hence the shell masses are respectively
\[
 1-p(T_k)^2,\qquad 1-p(T_k).
\]
This follows by integrating $2q\,dq$ and $dq$ through the transition;
it does not multiply delta currents. Beyond it the top-degree measures
vanish because both primitives are linear. On the shell $\phi_k=-k$,
so its absolute energy contributions are
\[
 k(1-p(T_k)^2),\qquad k(1-p(T_k)),
\]
and tend to zero by $1-p(T_k)\sim2/T_k$. Interior contributions converge
absolutely by the preceding integrability. The bounded psh potential
near $E$ contributes no mass to that analytic curve, also seen from
the zero radial boundary mass $p(-\infty)=0$.
Consequently the relative energies
\[
 E_\theta(\phi_k)=\frac13\sum_{j=0}^2
       \int_Y\phi_k(\theta+\ddc\phi_k)^j\wedge\theta^{2-j}
\]
have a finite limit and are uniformly bounded below. This proves
$\rho\in\mathcal E^1(Y,\theta)$ by the truncation definition
\cite[equation (2.1), Proposition 2.4 and Definition 2.5]{BBGZ}. By the uniqueness in the
problem this is its finite-energy solution. Equivalently, the full-mass
comparison principle of \cite[Corollary 2.3]{BEGZ} compares the two
exponential densities on each strict sublevel set and forces equality.

\section{Finite distance to the null curve}
Fix $v=v_0$ in an exceptional chart and put $h_0=1+|v_0|^2$.
The normal restriction is
\[
 \Omega\big|_{v=v_0}=\frac i{2\pi}
        \frac{p'(\log(|u|^2h_0))}{|u|^2}\,du\wedge d\bar u.
\]
Its coefficient is asymptotic to
$(3c\sqrt{h_0}/4)|u|^{-1}$ and diverges at every point of $E$.
Thus every positive smooth regular locus excludes $E$.
Along $u=r>0$ its squared real speed is bounded by $C/r$.
Its remaining length is at most $C\int_0^{r_0}r^{-1/2}\,dr<\infty$.
The divisor $D$ is also excluded from the smooth locus, since its normal
coefficient is asymptotic to $2/(|w|^2t^2)$ and diverges.

Let $U$ be any dense Zariski open positive regular locus and
$A=Y\setminus U$. On this compact surface $A$ has finitely many curve
and point components \cite[Theorem 7.8 and Corollary 7.10]{Adamus}, and contains $E$. Every other curve meets $E$
in only finitely many points. Choose a point of $E$ avoiding these
intersections and the point components. A sufficiently small coordinate
neighborhood has $A$ equal to $E$ there. The fixed-$v$ radial path
above therefore lies in $U$ for all sufficiently small positive $r$.
Its points approaching zero form an intrinsic Cauchy sequence, since
the tail lengths tend to zero. They have no limit in $U$: the smooth
positive metric induces the manifold topology, and their only limit
in $Y$ is the removed point of $E$. This proves incompleteness of every
such $U$.

The pole itself has a complete cusp end. On the punctured affine chart,
\[
 \Omega=p\ddc t+p'\,dt\wedge\dc t.
\]
Both terms are nonnegative and $p'\geq c/t^2$ for large $t$.
Length of any path to $D$ bounds below $c\int|dt|/t$, which diverges
as $t\to\infty$. The finite missing boundary is the disjoint null curve
$E$, not the prescribed pole. Theorem~\ref{thm:main} follows.

\section{Scope}
The smooth volume is chosen as part of the example. Exact identities
prove that it stays strictly positive across both ends; asymptotic
cancellation alone would not suffice. The construction retains the
actual divisorial equation, full mass, finite energy and minimal Lelong
numbers, and disproves completeness under mere nefness and bigness.
It does not assert failure or validity under the additional flow-compatible
data in \cite[Section 3.1]{DGL}; that branch remains open here.

\begin{thebibliography}{9}
\bibitem{Adamus} J.~Adamus, \emph{Complex Analytic Geometry},
lecture notes, revised 8 December 2015,
\href{https://www.math.uwo.ca/faculty/adamus/teaching/9607b2019/CAGI.pdf}{author's notes}.
\bibitem{BBGZ} R.~Berman, S.~Boucksom, V.~Guedj and A.~Zeriahi,
A variational approach to complex Monge--Amp\`ere equations,
\href{https://arxiv.org/abs/0907.4490}{arXiv:0907.4490v1}.
\bibitem{BEGZ} S.~Boucksom, P.~Eyssidieux, V.~Guedj and A.~Zeriahi,
Monge--Amp\`ere equations in big cohomology classes,
\emph{Acta Mathematica} 205 (2010), 199--262,
\href{https://arxiv.org/abs/0812.3674}{arXiv:0812.3674}.
\bibitem{BT} E.~Bedford and B.~A.~Taylor,
A new capacity for plurisubharmonic functions,
\emph{Acta Mathematica} 149 (1982), 1--40,
\href{https://www.math.stonybrook.edu/~ebedford/PapersForM655/BTNewCapacity.pdf}{author's copy}.
\bibitem{DGL} E.~Di Nezza, V.~Guedj and C.~H.~Lu,
Geometric smoothing by the K\"ahler--Ricci flow,
\href{https://arxiv.org/abs/2607.08325}{arXiv:2607.08325v1},
Theorem 3.4 and Question 3.5.
\bibitem{JP} M.~Jarnicki and P.~Pflug, \emph{Separately analytic functions},
manuscript, 11 November 2008,
\href{https://www.math.stonybrook.edu/~ebedford/PapersForM537/JarnickiPflugSepHolo081111.pdf}{source copy}.
\end{thebibliography}
\end{document}
